Data-driven weather models such as GraphCast, FourCastNet and Keisler's graph network are autoregressive: a network learns to map the atmospheric state at time $t$ to $t+\Delta$ and is applied repeatedly to produce a medium-range forecast~\cite{lam2022graphcast,pathak2022fourcastnet,keisler2022forecasting}. Neural general circulation models are trained through differentiable rollouts of many steps~\cite{kochkov2023neural}. Training on one step only leaves a mismatch with deployment: errors made at step $k$ become inputs at step $k{+}1$, and a model that is accurate one step ahead can drift off the attractor or blow up in long runs, a failure mode that has motivated noise injection, the pushforward trick, learned refinement, and post-hoc correction~\cite{stachenfeld2021learned,brandstetter2022message,lippe2023pde,pedersen2025thermalizer}. Rollout (multi-step) losses are a common remedy. The cost is a loss that is blind to the sharpness of the one-step map and a training cost that grows linearly with the rollout length.

How much of the benefit of rollout training is skill, how much is stability, and what does it do to climatological variance? Large systems bundle rollout fine-tuning with changes of architecture, data and resolution, so these effects are hard to separate. We study the question in the cheapest chaotic multiscale testbed that still has an unresolved-scale problem: the two-scale Lorenz-96 system (used as a testbed in the multi-scale learning studies of~\cite{chattopadhyay2019data,gagne2019machine}), where the emulator sees only the $36$ slow variables and the $360$ fast variables act as unobserved forcing.

We pose four falsifiable hypotheses, registered before running (Sec.~\ref{sec:setup}): (H1) multi-step training lowers medium-range RMSE; (H2) one-step training yields less stable free runs; (H3) multi-step training lowers climatological variance (blurring); (H4) the emulators beat persistence, climatology and a linear stencil in ACC lead time. We follow the WeatherBench convention of reporting RMSE and ACC against persistence and climatology as a function of lead time~\cite{rasp2020weatherbench,rasp2023weatherbench}.

\textbf{Contributions.} (i) A controlled sweep of the rollout length $K$ for one small CNN, five seeds per setting, with every system measured by the same evaluation code. (ii) A measured, mixed answer: a small, monotone medium-range gain from large $K$ (clear only for $K{=}8$), a one-step accuracy cost, and \emph{no} detectable effect on stability or variance. (iii) Two ablations that bound the result: input-noise training (harmful here, and the only intervention that visibly damps variance) and a $8\times$ smaller training set (no effect). The study is a small CPU study; Sec.~\ref{sec:limitations} lists what it cannot say.
